% Response to Reviewers -- round 2
% Manuscript: Influence of J2-plasticity on the effective crack resistance of a simplified
%             porous material in a numerical experiment
% Journal: Computational Science and Engineering (Springer Nature), collection KLAIM 2025
% Submission ID: 108f21b3-1ae0-46f3-912a-c62cdf4b765e
% Decision: major revision, 10 September 2026; deadline 1 November 2026 (extended)
%
% The quoted manuscript changes in the "Related changes" blocks are generated from
% manuscript_revised_template.tex by build_response_quotes.py (copy of 063/18_...):
%     pdflatex manuscript_revised_template.tex   (twice, so that the .aux is current)
%     python3 build_response_quotes.py manuscript_revised_template.tex manuscript_revised_template.aux response_to_reviewers_template.tex
%     pdflatex response_to_reviewers_template.tex (twice)
% Edit only the %@quote request lines and the text outside the %@generated blocks.
% Every \mycomment{TODO}{...} is an open item; the clean letter must contain none.

\documentclass[pdflatex,sn-mathphys-num]{sn-jnl}

\usepackage{graphicx}
\usepackage{amsmath,amssymb,amsfonts}
\usepackage{amsthm}
\usepackage{mathrsfs}
\usepackage{xcolor}
\usepackage{booktabs}
\usepackage[markup=underlined,authormarkup=none]{changes}
\usepackage{nicefrac}
\usepackage{xspace}

\DeclareRobustCommand{\mathdefault}[1]{#1}
\newcommand{\vMLabel}{[\textbf{vM}]\xspace}
\newcommand{\EqJcLabel}{[\textbf{Eq$\boldsymbol{J_c}$}]\xspace}
\newcommand{\EqSigmaStarLabel}{[\textbf{Eq$\boldsymbol{\sigma^*}$}]\xspace}
\newcommand{\changed}[1]{\begingroup\color{blue}#1\endgroup}
\newcommand{\mycomment}[2]{\textbf{\textcolor{red}{[#1: #2]}}}

\newenvironment{reviewercomment}
  {\par\noindent\textbf{Reviewer/editor comment.}\begin{quote}\itshape}
  {\end{quote}}
\newenvironment{response}
  {\par\noindent\textbf{Response.}\ }
  {\par}
\newenvironment{relatedchanges}
  {\par\noindent\textbf{Related changes in the manuscript.}\begin{itemize}}
  {\end{itemize}}
\newcommand{\msloc}[1]{\textit{#1:}\ }
\newcommand{\relhead}[1]{\item[]\vspace*{0.2\baselineskip}#1}

\title{Response to Reviewers (second revision)}

\author*[1]{\fnm{Alexander} \sur{Schl\"uter}}
\author[1]{\fnm{Ronjit} \sur{Medda}}
\author[1]{\fnm{Ralf} \sur{M\"uller}}
\affil*[1]{\orgdiv{Institut f\"ur Mechanik, FB Bau- und Umweltingenieurwissenschaften}, \orgname{TU Darmstadt},
  \orgaddress{\street{Franziska-Braun-Stra{\ss}e 7}, \city{Darmstadt}, \postcode{64287}, \country{Germany}}}

\begin{document}
\maketitle
\setlength{\parindent}{0pt}
\setlength{\parskip}{0.45\baselineskip}

\section*{Notes on this document}

The comments of the editor and the reviewers on the first revised version (v2.0, submitted 25 May 2026) are repeated
in italics and are followed by our response. Below each response, the related changes in the manuscript are listed as
bullets, one bullet per change, and quoted verbatim in the marked-up form of the revised manuscript: new text is
\added{underlined}, removed text is \deleted{struck out}, and changes inside equations are shown in \changed{blue}.
Each bullet starts with the location of the change in the revised manuscript. Equation, figure, table and reference
numbers in the quotes refer to the marked-up revised manuscript. A change that answers several comments is quoted once
and referred to elsewhere. All numerical results of the manuscript were recomputed for this revision (see the response
to Reviewer~1, point~3, and the section ``Additional corrections by the authors'').

\section*{Response to the Editor}

\begin{reviewercomment}
Ref: Submission ID 108f21b3-1ae0-46f3-912a-c62cdf4b765e. Decision: major revision. Deadline: 01 Oct 2026.

Dear Dr Schl\"uter, your manuscript, ``Influence of J2-plasticity on the effective crack resistance of a simplified
porous material in a numerical experiment'', has now been assessed. We invite you to revise your paper, carefully
addressing the comments from the reviewers and the editor. Please ensure the results are accurately reported, any
overstated conclusions are rewritten and the limitations of the work fully explained. When your revision is ready,
please submit the updated manuscript and a point-by-point response. This will help us move to a swift decision.

Kind regards, Mario Ohlberger, Editor, Computational Science and Engineering
\end{reviewercomment}

\begin{response}
We thank the editor and the three reviewers for their assessment and their constructive comments. We have revised the
manuscript accordingly and answer every comment point by point below. The main revisions are:
\begin{enumerate}
\item the introduction was rewritten (engineering problem $\to$ approaches and their limitations $\to$ previous contributions $\to$ open questions), the relation of the numerical experiment to homogenization is stated with Hill's definition of a representative volume element and with the rigorous results on the homogenization of surface energies, the term ``representative region'' was dropped, and the number of references was increased from 34 to 44;
\item \mycomment{TODO}{mechanism study: energy partition of the numerical experiment, dissipation per re-nucleation event, stress state in the ligament before re-nucleation (new Section~X / Figures~Y), sensitivity with respect to $\sigma_y$ and $H$, scaling argument for the linear elastic sets}
\item the square resp.\ cubic symmetry of the effective stiffness of the unit cells is quantified (Zener ratio) and the influence of the isotropic approximation of the embedding is tested by simulations with the full anisotropic stiffness\mycomment{TODO}{numbers from study S3};
\item all simulations recomputed with consistent parameters ($\beta = 0.2$, $\eta = 10^{-5}$ in two and three dimensions, linear elastic embedding); see ``Additional corrections by the authors''.
\end{enumerate}
\end{response}

%==================================================================================================
\section*{Reviewer 1}

\begin{reviewercomment}
Synopsis: The authors presents a study of the ``effective crack resistance'' of heterogeneous materials with pores where
the ``matrix'' material is governed by J2-plasticity and a phase-field fracture approach is used.

Evaluation: The topic is suited to CSE. The submitted version was significantly improved compared to the previous
version. Yet, certain central issues remain unresolved (to be specified below). Therefore, I recommend a major revision.
\end{reviewercomment}

\begin{response}
We thank the reviewer for the detailed second review and for acknowledging the improvements. We address the remaining
central issues below.
\end{response}

\subsection*{Major points}

\begin{reviewercomment}
1. The revision addressed many issues raised in the first review. However, several corrections concern matters of
notation, terminology, physical interpretation, and mathematical consistency that ideally should have been resolved
prior to (initial) submission (see minor points). Also, the level of scholarship feels uneven. My previous comments
forced the authors to introduce literature that arguably should have been there from the beginning. That gives the
impression that the literature review was assembled around a preferred narrative rather than developed from a broad
view of the field. This observation is supported by the low number of references.
\end{reviewercomment}

\begin{response}
We accept this criticism. The introduction was rewritten from a broader view of the field. It now distinguishes three classes of approaches to the crack resistance of heterogeneous materials and states the limitations of each: the rigorous homogenization of surface energies of Griffith type (Braides, Defranceschi and Vitali 1996; Cagnetti, Dal Maso, Scardia and Zeppieri 2019; the variational definition of Schneider and of Ernesti and Schneider; the assessment by Michel and Suquet), the analyses of crack-front trapping by perturbation and numerical methods (Gao and Rice 1989; Bower and Ortiz 1991; Lebihain, Leblond and Ponson 2020), and the numerical experiment of Hossain et al.\ with its recent developments. The literature on phase-field models coupled with plasticity was extended (Alessi, Marigo and Vidoli 2014; Ambati, Gerasimov and De Lorenzis 2015; Borden et al.\ 2016, in addition to Miehe et al., Kuhn et al.\ and Noll et al.), and Hill's definition of a representative volume element is cited where the term is discussed. In total, ten references were added (Hill 1963; Braides et al.\ 1996; Francfort and Marigo 1998; Cagnetti et al.\ 2019; Gao and Rice 1989; Bower and Ortiz 1991; Lebihain et al.\ 2020; Alessi et al.\ 2014; Borden et al.\ 2016; Ambati et al.\ 2015), so that the manuscript now cites 44 instead of 34 references, and the reference list is ordered by first citation. Independently of the reviewer's comments, we also went through the whole manuscript once more for notation, consistency and the correspondence between the text and the simulation code. The corrections resulting from this audit are listed in the section ``Additional corrections by the authors'' at the end of this letter.
\end{response}

\begin{relatedchanges}
\relhead{Introduction, rewritten (the new reference list entries are not quoted here, see the manuscript)}
%@quote cmd:Porous and other heterogeneous materials are increasingly used
%@quote cmd:For a heterogeneous material, the fracture resistance entering
%@quote cmd:For elastic and plastic properties, the step from the microstructure
%@quote cmd:Rigorous results exist for the homogenization of surface energies
%@quote cmd:A second class of approaches describes the interaction of a crack front
%@quote cmd:The third class is the numerical experiment of Hossain
%@quote cmd:A recent related development is the variable stiffness boundary condition of Hao et al.~\cite{hao_variable_2026}, which builds on the surfing boundary condition but translates a simple uniform remote displacement into a non-uniform load through an engineered stiffness gradient, thereby facilitating implementation in experiments and commercial finite-element software. The limitations
%@quote cmd:Phase-field models of fracture, see the review by Ambati
%@quote cmd:Our previous work considered only reversible elastic material behavior
%@quote cmd:The open questions addressed in this work are the following
%@quote cmd:A reliable determination of the fracture resistance parameters
%@quote cmd:The existence of an effective fracture toughness for a given material
%@quote cmd:assess essentially the same variational definition
%@quote cmd:In this work, we adopt the methodology introduced
%@quote cmd:The methodology is not used here as a proof
%@generated-begin
\item \msloc{Introduction}
\replaced{Porous and other heterogeneous materials are increasingly used in load-bearing components. Additive manufacturing processes such as wire arc additive manufacturing (WAAM) allow materials to be combined and microstructures to be designed through the process parameters, as demonstrated by Hauser et al.~[2], and pores are a typical defect of these processes. The strength assessment of such defect-rich materials is often based on fracture mechanics, which requires a fracture resistance parameter of the material. Relevant quantities are the critical stress intensity factor $K_{Ic}$ of Irwin's criterion $K_I = K_{Ic}$, the critical energy release rate $G_c$ of Griffith's criterion $G = G_c$ and, more generally, a critical value $J_c$ of the J-integral $J$ in the criterion}{While homogenization techniques for elastic and plastic material properties are well established,
the homogenization of fracture toughness remains a comparatively young field of research.
Relevant quantities include the critical stress intensity factor \( K_{Ic} \),
appearing in Irwin’s criterion \( K_I = K_{Ic} \), and the critical energy release rate \( G_c \),
used in Griffith’s criterion \( G = G_c \).
Similarly,  more general fracture formulations based on the J-integral $J$ can be established, i.e.}
\item \msloc{Introduction}
\added{For a heterogeneous material, the fracture resistance entering eq.~(1) is not a property of the bulk material alone but depends on the microstructure. How it can be determined for a porous material with a ductile bulk material is the subject of this work.}
\item \msloc{Introduction}
\deleted{A reliable determination of the fracture resistance parameters which are properties of the
underlying material and microstructure is of significant
practical interest, particularly in light of the increasing use of
additive manufacturing technologies. For example, wire arc additive manufacturing (WAAM) enables the flexible combination of materials
and targeted microstructural design through the control of process parameters,
as demonstrated by Hauser et al.~[2]. Fracture-mechanics-based criteria are often the appropriate tools
for strength assessment in these defect-rich materials.
However, their application requires material- and microstructure-dependent
fracture resistance parameters as mentioned above.}
\item \msloc{Introduction}
\added{For elastic and plastic properties, the step from the microstructure to a macroscopic material model is provided by homogenization. Its central concept is the representative volume element, defined by Hill~[4] as the smallest sample of a heterogeneous material that can be treated as a statistically uniform continuum. The effective properties of such a sample are independent of the boundary conditions applied to it, see also Nemat-Nasser and Hori~[5]. Fracture does not fit into this framework in a direct way. A crack is a localized and evolving discontinuity, its growth is a non-stationary process that may be unstable, and the separation of length scales that underlies the notion of a representative volume element is lost once the crack interacts with individual microstructural features. For this reason, the term \textsl{fracture toughness} is used in this work only for a macroscopic material property in the sense of eq.~(1), whereas the value obtained from a numerical experiment under specific conditions, i.e., the maximum energetic driving force required to force crack advance through the considered microstructure, is called the \textsl{effective crack resistance}. Three classes of approaches to the latter quantity can be distinguished in the literature.}
\item \msloc{Introduction}
\added{Rigorous results exist for the homogenization of surface energies of Griffith type. Braides, Defranceschi and Vitali~[6] proved that free-discontinuity functionals with a periodically oscillating surface energy, as they appear in the variational formulation of brittle fracture by Francfort and Marigo~[7] and Bourdin et al.~[8], $\Gamma$-converge to a functional with a homogenized and, in general, anisotropic surface energy. Cagnetti, Dal Maso, Scardia and Zeppieri~[9] extended this result to stationary random microstructures. Schneider~[10] and Ernesti and Schneider~[11, 12] use the resulting cell formula as a variational definition of an effective crack energy and compute it with FFT-based solvers, also for anisotropic fracture resistance. Michel and Suquet~[13] assess the merits and limitations of this definition. They relate it to limit-load problems and show that it is a lower bound of a more physically motivated effective fracture energy, namely the energy dissipated along the actual crack path through a unit cell. The homogenized surface energy is the energy of the minimal crack surface weighted by the local toughness. It does not account for the elastic interaction of the crack with the microstructure, for crack trapping and the unstable crack jumps that follow re-nucleation, or for inelastic bulk dissipation, which are the mechanisms that govern toughening in many materials, see the discussion by Lebihain~[14].}
\item \msloc{Introduction}
\added{A second class of approaches describes the interaction of a crack front with the microstructure directly. Gao and Rice~[15] and Bower and Ortiz~[16] analyzed the trapping of a crack front by arrays of tough obstacles with perturbation methods and with numerical methods for three-dimensional cracks, and Lebihain, Leblond and Ponson~[17] derived the effective toughness of periodic heterogeneous materials including out-of-plane excursions of the crack. These approaches are restricted to linear elastic bulk behavior and to small contrasts of the toughness or small perturbations of a planar crack.}
\item \msloc{Introduction}
\added{The third class is the numerical experiment of Hossain et al.~[1]. A region of the material that resolves the microstructure is subjected to so-called surfing boundary conditions, which enforce complete separation at a prescribed macroscopic crack speed, and the maximum of the energetic driving force recorded at the boundary is taken as the effective crack resistance. Since the crack evolution inside the region is computed with a phase-field model, crack trapping, re-nucleation, branching and unstable crack jumps are included without restriction. The method was developed further by Hsueh et al.~[18] for layered materials and by Brach et al.~[19] for anisotropic effective crack resistance.}
\item \msloc{Introduction}
\added{A recent related development is the variable stiffness boundary condition of Hao et al.~[24], which builds on the surfing boundary condition but translates a simple uniform remote displacement into a non-uniform load through an engineered stiffness gradient, thereby facilitating implementation in experiments and commercial finite-element software. The limitations of this third class are of a different nature than those of the first two. The recorded value depends on the loading of the experiment and on the size and the content of the resolved region, and the representativity of this region in the sense of Hill is not guaranteed but has to be assessed for every microstructure. The result is therefore an operational quantity of the experiment and not a homogenized material property. With few exceptions, the contributions cited above are limited to linear elastic material behavior and to two-dimensional problems.}
\item \msloc{Introduction}
\added{Phase-field models of fracture, see the review by Ambati et al.~[25], have been coupled with plasticity in several ways. Gradient damage models coupled with plasticity and their relation to cohesive cracks were analyzed by Alessi, Marigo and Vidoli~[26]. Variational formulations at finite strains were proposed by Miehe, Aldakheel and co-workers~[27, 28] and by Borden et al.~[29], and small-strain formulations in which the plastic energy is degraded together with the elastic energy were proposed by Ambati, Gerasimov and De Lorenzis~[30] and by Kuhn et al.~[31] and Noll et al.~[32]. In the latter formulation, a cubic degradation function allows the yield stress and the hardening modulus to be prescribed directly, and it is the formulation used in this work.}
\item \msloc{Introduction}
\deleted{The existence of an effective fracture toughness for a given material
and the question of how it should be defined remain subjects
of active debate, as discussed by Lebihain~[14],
with several approaches proposed in the literature.
In contrast to the effective elastic or plastic properties obtained in classical micromechanics, it is presently not established in general that fracture in a heterogeneous material gives rise to a homogenized macroscopic material model with a well-defined fracture toughness obtained in an RVE limit.
To avoid conflating two slightly different concepts, the term effective crack resistance is used in this work for the value obtained from a numerical experiment under specific less general conditions, i.e., the maximum energetic driving force required to force crack advance through the considered microstructure. In contrast, fracture toughness is reserved for a generalized material quantity in the sense of a macroscopic fracture property. Whether such a fracture toughness exists for heterogeneous materials and under which circumstances both concepts might coincide remains an open question of research.
Schneider~[10] suggests a variational definition that
relies on rigorous mathematical results for the homogenization of
the energy functional occurring in Bourdin's popular variational formulation of brittle fracture~[8].
This idea is further pursued by
Ernesti and Schneider~[11, 12] for anisotropic fracture resistance.}
\item \msloc{Introduction}
\deleted{Michel and Suquet~[13]
assess essentially the same variational definition and discuss its merits and limitations. In particular, they relate the variational surface energy to limit-load problems and show that it provides a lower bound to a more physically motivated effective fracture energy based on the energy dissipated along the actual crack path through a unit cell.}
\item \msloc{Introduction}
\deleted{In this work, we adopt the methodology introduced
by Hossain et al.~[1] and further
developed and applied by Hsueh et al.~[18] for layered materials,
as well as by Brach et al.~[19] for anisotropic effective crack resistance,
to determine the effective crack resistance of heterogeneous materials.
All of these contributions were limited to linear-elastic material behavior and two-dimensional problems.}
\item \msloc{Introduction}
\deleted{The methodology is not used here as a proof of a homogenized fracture toughness in the mathematical sense. Rather, it is interpreted as a numerical counterpart of a fracture experiment.
In this experiment-like setting,
a representative material region containing
 microstructural features is subjected
 to so-called surfing boundary conditions that enforce complete separation, and the energetic driving force required for this process is recorded.
  A recent related development is the variable stiffness boundary condition of Hao et al.~[24], which builds on the surfing boundary condition but translates a simple uniform remote displacement into a non-uniform load through an engineered stiffness gradient, thereby facilitating implementation in experiments and commercial finite-element software.
The resulting microstructural crack evolution is
then simulated using a phase-field model.
See, for example, the review by Ambati et al.~[25] for a comprehensive overview of phase-field models for fracture.
For ductile materials, phase-field fracture has also been coupled to plasticity in variational finite-strain formulations, for example by Miehe, Aldakheel and co-workers~[27, 28].}
\item \msloc{Introduction}
\replaced{Our previous work considered only reversible elastic material behavior. The Ramberg--Osgood law is non-linear and coincides with von Mises plasticity under monotonic loading, but it does not capture unloading. Here, the effect of true energy dissipation due to plasticity is studied with a von Mises plasticity model.}{In contrast to our previous work, which only considered reversible elastic material behavior,
albeit non-linear and consistent with von Mises plasticity under monotonic loading,
but which does not accurately capture unloading behavior,
we here also study the effect of true energy dissipation due to plasticity in ductile materials.}
\item \msloc{Introduction}
\replaced{The open questions addressed in this work are the following. First, how does plastic dissipation change the effective crack resistance obtained in the numerical experiment, and by which mechanisms? Compared with a physical experiment, the simulation gives access to the complete energy balance of the specimen and to the stress state at the location of crack re-nucleation. Both are used here to separate the contributions of crack nucleation, plastic dissipation and crack growth, and the sensitivity of the result with respect to the yield stress and the hardening modulus is examined. Second, does the two-dimensional idealization of cylindrical pores carry over to spherical pores in three dimensions? Third, by how much does the result depend on the initial crack and on the size of the resolved region? The last point does not replace a representativity study, but it quantifies the dependence for the microstructure considered.}{In addition, we extend the approach of Hossain et al.~[1] to take into account the three-dimensional
character of real porous materials with spherical pores.}
%@generated-end
\end{relatedchanges}

\begin{reviewercomment}
2. Although the authors did change the manuscript to some extent, the understanding of homogenization still seems
shaky. I still do not get the feeling that the authors have a clear conceptual framework connecting: homogenization,
effective properties, fracture-process representations, RVE limits and the Hossain numerical experiment. For instance,
the very first sentence of the introduction reads ``While homogenization techniques for elastic and plastic material
properties are well established, the homogenization of fracture toughness remains a comparatively young field of
research.'' There is no homogenization of fracture toughness!!!
\begin{itemize}
\item Homogenization is a process with well-defined input and output. Once you homogenize milk, it is homogeneous.
Periodic/stochastic homogenization in linear elasticity works similarly. This is not true for fracture toughness.
\item ``Representative'' is a quality that is earned and that can be assessed. For instance, Hill defined an RVE to be
``the smallest material sample of a heterogeneous mixture that can be treated as a statistically uniform continuum''.
The authors use the wording ``representative region'' although it is just a sample/snapshot and it is not clear that it
is representative.
\item The authors down-toned their claims of representativity of the strategy. However, the manuscript still reads as
if this property would be automatic and the current work just did not study it. I firmly disagree! It remains entirely
unclear whether the Hossain experiment would withstand a rigorous representativity study.
\item A material property is a property of the material. In particular, it does not depend on specific choices like the
boundary conditions. Quite generally, it is not clear that a microstructured material at the macroscopic scale can be
described by such a coupled phase-field-plasticity model and what the parameters would look like.
\end{itemize}
\end{reviewercomment}

\begin{response}
We agree with the reviewer on all four points and have changed the conceptual framing of the manuscript accordingly.
\begin{itemize}
\item The first sentence and the whole introduction were rewritten. The manuscript no longer speaks of a ``homogenization of fracture toughness''. It now states what homogenization provides for elastic and plastic properties (Hill's representative volume element, whose effective properties are independent of the boundary conditions), why fracture does not fit into this framework directly (localized, evolving, possibly unstable process; loss of the separation of length scales), and which rigorous results exist nevertheless for surface energies of Griffith type (Braides et al.\ 1996, Cagnetti et al.\ 2019, the cell formula used by Schneider and assessed by Michel and Suquet). The quantity determined in this work is consistently called the effective crack resistance of the numerical experiment and is explicitly distinguished from a fracture toughness in the sense of a material property.
\item ``Representative'' is no longer used for the sample. The region is called the resolved porous region (or resolved region), Hill's definition of a representative volume element is quoted in Section~2, and it is stated that the representativity of the resolved region in this sense is not established. The term ``representative'' was also removed from the figure caption of Fig.~1, from the conclusion and, in its colloquial sense (``representative of the main trend''), from Section~4.
\item The manuscript no longer suggests that representativity would be automatic. Section~2 now states that the independence of $J_{x,c}^{\text{eff}}$ of the initial crack and of the size of the resolved region cannot be taken for granted and has to be tested for every microstructure, that only one crack-growth direction is considered, and that a representativity study in the sense of Hill (statistics over microstructure realizations, sample sizes and loading conditions) is outside the scope of this work. What is done instead is a direct probe of the two dependencies that can be tested within the present setup: the length of the initial crack and the number of resolved pore rows are varied in the new Section~\mycomment{TODO}{4.6 (study S5)}, and the resulting change of $J_{x,c}^{\text{eff}}$ is reported there.\mycomment{TODO}{numbers from study S5}
\item We do not claim that the microstructured material is described by the coupled phase-field--plasticity model at the macroscopic scale. The macroscopic use of $J_{x,c}^{\text{eff}}$ is restricted to a crack-growth criterion of the type of eq.~(1) under the loading of the experiment, and the dependence of the recorded value on the loading of the experiment is stated as a limitation of the approach in the introduction and in Section~2.
\end{itemize}
\end{response}

\begin{relatedchanges}
\relhead{Section~2, framing of the numerical experiment (the rewritten introduction is quoted under major point~1)}
%@quote cmd:A Numerical Experiment for the Crack Resistance of a Resolved Porous Region
%@quote cmd:that resolves the}{containing representative}
%@quote cmd:finite region of the material
%@quote cmd:This region is referred to as the resolved porous region
%@quote cmd:This definition does not imply that
%@quote cmd:cannot be taken for granted and has to be tested
%@quote cmd:resolved porous region}{representative
%@generated-begin
\item \msloc{Section~2}
\replaced{A Numerical Experiment for the Crack Resistance of a Resolved Porous Region}{A Numerical Experiment to Determine the Effective Crack Resistance of Materials with Complex Microstructure}
\item \msloc{Caption of Fig.~1}
\replaced{that resolves the}{containing representative}
\item \msloc{Section~2}
\replaced{finite region of the material}{representative region}
\item \msloc{Section~2}
\replaced{This region is referred to as the resolved porous region, or simply the resolved region, in the following. The term representative volume element is deliberately avoided. Hill~[4] defines a representative volume element as the smallest sample of a heterogeneous material that can be treated as a statistically uniform continuum, so that its effective properties are independent of the applied boundary conditions. The resolved region considered here is a sample of the microstructure whose representativity in this sense is not established, and the separation of length scales underlying Hill's definition is lost once a crack of the size of the region develops.}{The term representative volume element is deliberately avoided because the region no longer fulfills classical representativity requirements once fracture occurs, e.g., the separation of length scales.}
\item \msloc{Section~2}
\replaced{This definition does not imply that $J_{x,c}^{\text{eff}}$ becomes independent of the size of the resolved region and of the position of the initial crack when the region is enlarged. Whether this is the case for a given class of microstructures and loading paths is a question of representativity that is not settled by the definition.}{This definition should not be confused with a proof that $J_{x,c}^{\text{eff}}$ converges to a material property under enlargement of the representative region. Establishing such convergence, if it exists for a given class of microstructures and loading paths, would require a separate study.}
\item \msloc{Section~2}
\replaced{Whether $J_{x,c}^{\text{eff}}$ is independent of the initial microscopic crack configuration and of the size of the resolved region cannot be taken for granted and has to be tested for every microstructure. In the present work, only one macroscopic crack-growth direction is considered. The influence of the length of the initial crack and of the number of resolved pore rows on $J_{x,c}^{\text{eff}}$ is probed in Section~4.6. A representativity study in the sense of Hill~[4], i.e., a statistical assessment over microstructure realizations, sample sizes and loading conditions, is outside the scope of this work.}{Independence from the initial microscopic crack configuration and adequate representation of fracture mechanisms can be ensured by selecting a sufficiently large representative region that contains an appropriate number of microstructural features.
This presumes a sufficiently homogeneous spatial distribution of microstructural features, analogous to assumptions made for representative volume elements in micromechanics.
In the present work, this requirement is not fully verified: only one prescribed macroscopic crack-growth direction and one initial crack location are considered, and no convergence study with respect to the size of the resolved porous region is performed.}
\item \msloc{Conclusion}
\replaced{resolved porous region}{representative
microstructure}
%@generated-end
\end{relatedchanges}
\begin{relatedchanges}
\relhead{\mycomment{TODO}{quotes: new Section 4.6 with the S5 results}}
\end{relatedchanges}

\begin{reviewercomment}
3. Results are often interpreted more strongly than justified, but there is no real mechanism study. No energy
partition, no sensitivity study, no reduced model, no scaling argument, no asymptotic analysis, no attempt to isolate
competing effects. So many conclusions are essentially post hoc physical interpretations of numerical observations.
That is not necessarily wrong, but it is not especially rigorous.
\end{reviewercomment}

\begin{response}
\mycomment{TODO}{answer: (a) energy partition of every simulation (exact balance $W = E_{\mathrm{el}} + E_{\mathrm{pl}} + E_{\mathrm{fr}} + D_{\mathrm{visc}}$, per re-nucleation event; the energy released in the unstable ligament failure is dissipated by the viscous regularisation and is the quantity that distinguishes the material laws); (b) stress state and plastic zone in the ligament immediately before re-nucleation (event snapshots) for the two regimes $w_s \lesssim 0.6L$ and $w_s \gtrsim 2L$; (c) sensitivity study $\sigma_y \in \{0.75, 1.5\}$, $H \in \{0.11, 0.44\}\mu_0$ at $w_s = 0.375L$ and $1.0L$ (study S4); (d) scaling argument: for the rate-independent linear elastic model $J_{x,c}^{\text{eff}} \propto J_c$ exactly, confirmed by \EqJcLabel vs.\ \EqSigmaStarLabel; (e) reduced one-dimensional ligament model for the nucleation regime (Sec.~3.6 extended). Numbers from 075/076 results.}
\end{response}

\begin{relatedchanges}
\relhead{\mycomment{TODO}{quotes: new mechanism section, new figures (energy partition vs.\ $w_s$, event energies, ligament stress), sensitivity paragraph, scaling paragraph}}
\end{relatedchanges}

\begin{reviewercomment}
4. ``The representative region is embedded within a surrounding matrix characterized by effective elastic properties,
such as the effective Young's modulus $E^{\mathrm{eff}}$ and Poisson's ratio $\nu^{\mathrm{eff}}$ in the isotropic
case.'' $\to$ This assumption assumes isotropy of the effective elastic behavior. Interestingly, \emph{none} of the
considered examples in the manuscript leads to isotropic effective elastic properties. Both the 2d and the 3d porous
microstructure are \emph{cubic}, not isotropic. If isotropic effective elastic properties are used although the
microstructures possess only cubic symmetry, it is unclear whether the embedding problem is posed consistently. The
authors should discuss the implications of this approximation and justify why it does not affect the reported
conclusions. Also, embedding the heterogeneous material in a homogeneous material leads to additional discontinuities
and notch effects due to the sharp corners. This issue is not discussed.
\end{reviewercomment}

\begin{response}
The reviewer is right. The unit cells of both microstructures have square (two-dimensional) resp.\ cubic (three-dimensional) symmetry, and the previous versions assigned an isotropic stiffness to the embedding without saying so or quantifying the approximation. The revised manuscript states the symmetry in Section~2, reports the full effective stiffness of every unit cell in Sections~4.1 and 5.1 ($C_{11}$, $C_{12}$, $C_{33}$ resp.\ $C_{44}$, and the Zener ratio $A = 2C_{33}/(C_{11}-C_{12})$ resp.\ $2C_{44}/(C_{11}-C_{12})$, Tables~\mycomment{TODO}{X and Y}), and explains which constants the isotropic approximation uses. For the two-dimensional unit cells, the Zener ratio is $0.89$ for the smallest wall width $w_s = 0.375L$, $0.96$ for $w_s = L$ and $1.00$ for $w_s = 4L$, i.e., the anisotropy is moderate and vanishes for widely spaced pores. The influence of the approximation on the result is tested directly: the simulations of Section~\mycomment{TODO}{4.5 (study S3)} assign the full square-symmetric (cubic) stiffness to the embedding for $w_s = 0.375L$, $1.0L$ and $4.0L$ in two dimensions and for $w_s = 0.375L$ and $1.0L$ in three dimensions.\mycomment{TODO}{result: $J_{x,c}^{\text{eff}}$ changes by $y$\,\% (study S3), which justifies the approximation for the conclusions drawn.}

Regarding the corners: the transition from the resolved porous region to the homogeneous embedding is a material discontinuity and the corners of the resolved region are stress concentrators. This is now stated in Section~2. The corners are located at least one cell width $w_c$ away from the crack plane and at least one cell width ahead of the first and behind the last pore, the embedding is linear elastic (the yield stress is not assigned to it, which is now stated explicitly in Section~4.1), and the J-integral is evaluated on the outer boundary of the embedding. The local stress concentration at the corners therefore does not enter the recorded driving force, and it did not lead to crack nucleation at a corner in any of the simulations.\mycomment{TODO}{quantify the corner stress from the field output of study S1 and state it in Section~4.5}
\end{response}

\begin{relatedchanges}
\relhead{Sections~2, 4.1 and 5.1 (embedding, symmetry of the effective stiffness)}
%@quote cmd:homogeneous matrix with the effective elastic stiffness of a unit cell
%@quote cmd:The transition from the resolved region to the homogeneous embedding
%@quote cmd:The unit cell has square symmetry
%@quote cmd:The embedding is therefore modelled as linear elastic
%@quote cmd:The unit cell has cubic symmetry
%@generated-begin
\item \msloc{Section~2}
\replaced{homogeneous matrix with the effective elastic stiffness of a unit cell of the microstructure. For the periodic arrangements of circular or spherical pores considered in this work, this stiffness has square or cubic symmetry. In the simulations it is approximated by an isotropic stiffness with the effective Young's modulus $E^{\text{eff}}$ and Poisson's ratio $\nu^{\text{eff}}$. The degree of anisotropy of the unit cells and the influence of this approximation on the effective crack resistance are quantified in Sections~4.1 and~5.1.}{matrix characterized by effective elastic properties,
such as the effective Young's modulus $E^{\text{eff}}$ and Poisson's ratio $\nu^{\text{eff}}$ in the isotropic case.}
\item \msloc{Section~2}
\added{The transition from the resolved region to the homogeneous embedding is a material discontinuity, and the corners of the resolved region act as geometric stress concentrators. In the simulations of Sections~4 and~5, the corners are located at least one cell width $w_c$ away from the crack plane and the embedding remains linear elastic, so that the J-integral evaluated on the outer boundary is not affected by the local stress concentration at the corners. The influence of the corners on the crack pattern is discussed together with the results.}
\item \msloc{Section~4.1}
\added{The unit cell has square symmetry, so that $\left< \underline{C} \right>_{11} = \left< \underline{C} \right>_{22}$ and the effective stiffness is characterized by the three independent constants $C_{11} = \left< \underline{C} \right>_{11}$, $C_{12} = \left< \underline{C} \right>_{12}$ and $C_{33} = \left< \underline{C} \right>_{33}$. The isotropic approximation of eq.~(36) uses $C_{12}$ and $C_{33}$ and replaces $C_{11}$ by $C_{12} + 2C_{33}$. The deviation from isotropy is measured by the Zener ratio $A = 2C_{33}/(C_{11} - C_{12})$, which equals one for an isotropic material. Table~1 lists the effective constants and the Zener ratio for all wall widths considered. The anisotropy is largest for the smallest wall width and vanishes for widely spaced pores. The influence of the isotropic approximation on the effective crack resistance is tested in Section~4.5 by simulations in which the embedding is assigned the full square-symmetric stiffness.}
\item \msloc{Section~4.1}
\added{The embedding is therefore modelled as linear elastic, i.e., no yielding is admitted outside the resolved porous region.}
\item \msloc{Section~5.1}
\added{The unit cell has cubic symmetry with the three independent constants $C_{11}$, $C_{12}$ and $C_{44}$ and the Zener ratio $A = 2C_{44}/(C_{11}-C_{12})$, which are listed in Table~2.}
%@generated-end
\end{relatedchanges}
\begin{relatedchanges}
\relhead{\mycomment{TODO}{quotes: Tables of the effective constants, new Section 4.5 with the S3 results}}
\end{relatedchanges}

\subsection*{Minor points}

\begin{reviewercomment}
0. Avoid the use of semicolons. Use (different, clear, legible) sentences.
\end{reviewercomment}
\begin{response}
All semicolons in running text were replaced by separate sentences (ten occurrences). Three examples are quoted below.
\end{response}
\begin{relatedchanges}
\relhead{Sections~2, 4.1 and 5.2 (examples)}
%@quote cmd:. Hence,}{; hence,}
%@quote cmd:. Several}{; several}
%@quote cmd:. It}{; it}
%@generated-begin
\item \msloc{Section~2}
\replaced{. Hence,}{; hence,}
\item \msloc{Section~4.1}
\replaced{. Several}{; several}
\item \msloc{Section~5.2}
\replaced{. It}{; it}
%@generated-end
\end{relatedchanges}

\begin{reviewercomment}
1. The introduction feels wrong. In the first line, it is claimed that the fracture toughness can be homogenized (which
is not clear imho). The subsequent discussion appears unordered. The introduction should clearly introduce the overall
problem (for heterogeneous materials, fracture toughness is required for engineering design), different approaches to
tackle it, together with the limitations (ad-hoc models, bounds, Hossain's approaches). Then, it should discuss previous
contributions to the area and open questions.
\end{reviewercomment}
\begin{response}
The introduction was restructured in exactly this order: the engineering problem (fracture resistance parameters are needed for heterogeneous, defect-rich materials), the approaches to the problem together with their limitations (rigorous homogenization of surface energies, crack-front perturbation analyses, the numerical experiment of Hossain et al.), previous contributions, and the open questions addressed in this work. The claim of the first sentence was removed. The rewritten introduction is quoted under major point~1.
\end{response}

\begin{reviewercomment}
2. The sentence ``. In contrast to \dots'' is way too long and incomprehensible.
\end{reviewercomment}
\begin{response}
The sentence was split into three sentences.
\end{response}
\begin{relatedchanges}
\relhead{Introduction}
%@quote cmd:Our previous work considered only reversible elastic material behavior
%@generated-begin
\item \msloc{Introduction}
\replaced{Our previous work considered only reversible elastic material behavior. The Ramberg--Osgood law is non-linear and coincides with von Mises plasticity under monotonic loading, but it does not capture unloading. Here, the effect of true energy dissipation due to plasticity is studied with a von Mises plasticity model.}{In contrast to our previous work, which only considered reversible elastic material behavior,
albeit non-linear and consistent with von Mises plasticity under monotonic loading,
but which does not accurately capture unloading behavior,
we here also study the effect of true energy dissipation due to plasticity in ductile materials.}
%@generated-end
\end{relatedchanges}

\begin{reviewercomment}
3. The title of section ``2 A Numerical Experiment to Determine the Effective Crack Resistance of Materials with Complex
Microstructure'' suggests that the effective crack resistance is a material property.
\end{reviewercomment}
\begin{response}
The title of Section~2 was changed so that it refers to the resolved porous region and not to a material.
\end{response}
\begin{relatedchanges}
\relhead{Title of Section~2}
%@quote cmd:A Numerical Experiment for the Crack Resistance of a Resolved Porous Region
%@generated-begin
\item \msloc{Section~2}
\replaced{A Numerical Experiment for the Crack Resistance of a Resolved Porous Region}{A Numerical Experiment to Determine the Effective Crack Resistance of Materials with Complex Microstructure}
%@generated-end
\end{relatedchanges}

\begin{reviewercomment}
4. equations are not words, e.g., ``by (1)'' $\to$ ``by eq. (1)''
\end{reviewercomment}
\begin{response}
All references to equations are now written as ``eq.~(n)'' (eleven places were changed, one of which is quoted below).
\end{response}
\begin{relatedchanges}
\relhead{Section~2 (example)}
%@quote cmd:eq.~@1
%@generated-begin
\item \msloc{Introduction}
\added{For a heterogeneous material, the fracture resistance entering eq.~(1) is not a property of the bulk material alone but depends on the microstructure. How it can be determined for a porous material with a ductile bulk material is the subject of this work.}
%@generated-end
\end{relatedchanges}

\begin{reviewercomment}
5. ``representative region'' (see major point 2)
\end{reviewercomment}
\begin{response}
The term was removed everywhere, see major point~2. The region is called the resolved porous region.
\end{response}
\begin{relatedchanges}
\relhead{Section~2 and conclusion}
%@quote cmdall:resolved}{representative}
%@generated-begin
\item \msloc{Section~2}
\replaced{resolved}{representative}
\item \msloc{Section~2}
\replaced{resolved}{representative}
\item \msloc{Section~2}
\replaced{resolved}{representative}
\item \msloc{Section~2}
\replaced{resolved}{representative}
\item \msloc{Section~2}
\replaced{resolved}{representative}
\item \msloc{Section~2}
\replaced{resolved}{representative}
%@generated-end
\end{relatedchanges}

\begin{reviewercomment}
6. ``Furthermore, the simulation must be based on physically consistent principles.'' The statement does not appear to
convey actionable information and should either be clarified or removed.
\end{reviewercomment}
\begin{response}
The sentence was removed.
\end{response}
\begin{relatedchanges}
\relhead{Section~2}
%@quote cmd:Furthermore, the simulation must be based on physically consistent principles.
%@generated-begin
\item \msloc{Section~2}
\deleted{Furthermore, the simulation must be based on physically consistent principles.}
%@generated-end
\end{relatedchanges}

\begin{reviewercomment}
7. ``Only small strain assumptions are considered in this work.'' $\to$ ``Only small strains are considered in this work.''
\end{reviewercomment}
\begin{response}
Done.
\end{response}
\begin{relatedchanges}
\relhead{Section~2}
%@quote cmd: assumption@1
%@generated-begin
\item \msloc{Section~2}
\deleted{ assumption}
%@generated-end
\end{relatedchanges}

\begin{reviewercomment}
8. ``The recorded value is \dots'' $\to$ not a sentence. Rewrite for clarity.
\end{reviewercomment}
\begin{response}
The passage was rewritten as three sentences.
\end{response}
\begin{relatedchanges}
\relhead{Section~2}
%@quote cmd:The recorded value depends on the crack pattern
%@generated-begin
\item \msloc{Section~2}
\replaced{The recorded value depends on the crack pattern inside the resolved microstructure. Following the observation reported by Hossain et al.~[1], the homogeneous surrounding material is assumed not to change the value of the boundary J-integral. Under this assumption, the J-integral evaluated at the outer boundary of the embedding equals the macroscopic driving force acting on the resolved region.}{The recorded value is, of course, affected by the crack pattern inside the explicitly resolved microstructure; however, following the observation reported by Hossain et al.~[1], the additional homogeneous surrounding material is assumed not to change the value of the boundary J-Integral, so that evaluating it at the boundary of the resolved microstructure or at the boundary of an embedding region with representative effective properties gives the same macroscopic driving force.}
%@generated-end
\end{relatedchanges}

\begin{reviewercomment}
9. ``The index c denotes the critical crack-resistance value obtained from this maximum,'' $\to$ ``The index c indicates
that the crack-resistance value is obtained as a maximum, i.e., is \emph{critical}''
\end{reviewercomment}
\begin{response}
The suggested wording was adopted.
\end{response}
\begin{relatedchanges}
\relhead{Section~2}
%@quote cmd:indicates that the crack-resistance value is obtained as a maximum
%@generated-begin
\item \msloc{Section~2}
\replaced{indicates that the crack-resistance value is obtained as a maximum, i.e., is \textsl{critical}}{denotes the critical crack-resistance value obtained from this maximum}
%@generated-end
\end{relatedchanges}

\begin{reviewercomment}
10. ``Establishing such convergence, if it exists for a given class of microstructures and loading paths, would require
a separate study.'' Convergence is a mathematical property, not an empirical observation. Please clarify what is meant
here. $\to$ see major item 2
\end{reviewercomment}
\begin{response}
The word ``convergence'' was inappropriate. The sentence now states what is meant, namely whether $J_{x,c}^{\text{eff}}$ becomes independent of the size of the resolved region and of the position of the initial crack when the region is enlarged, and that this is a question of representativity which the definition does not settle. See also major point~2.
\end{response}
\begin{relatedchanges}
\relhead{Section~2}
%@quote cmd:This definition does not imply that
%@generated-begin
\item \msloc{Section~2}
\replaced{This definition does not imply that $J_{x,c}^{\text{eff}}$ becomes independent of the size of the resolved region and of the position of the initial crack when the region is enlarged. Whether this is the case for a given class of microstructures and loading paths is a question of representativity that is not settled by the definition.}{This definition should not be confused with a proof that $J_{x,c}^{\text{eff}}$ converges to a material property under enlargement of the representative region. Establishing such convergence, if it exists for a given class of microstructures and loading paths, would require a separate study.}
%@generated-end
\end{relatedchanges}

\begin{reviewercomment}
11. ``Independence from the initial microscopic crack configuration and adequate representation of fracture mechanisms
can be ensured by selecting a sufficiently large representative region that contains an appropriate number of
microstructural features.'' Is that so? $\to$ see major item 2
\end{reviewercomment}
\begin{response}
No, it is not. The sentence was removed and replaced by the statement that this independence cannot be taken for granted, by a description of what is tested in this work and of what is not, see major point~2.
\end{response}
\begin{relatedchanges}
\relhead{Section~2}
%@quote cmd:cannot be taken for granted and has to be tested
%@generated-begin
\item \msloc{Section~2}
\replaced{Whether $J_{x,c}^{\text{eff}}$ is independent of the initial microscopic crack configuration and of the size of the resolved region cannot be taken for granted and has to be tested for every microstructure. In the present work, only one macroscopic crack-growth direction is considered. The influence of the length of the initial crack and of the number of resolved pore rows on $J_{x,c}^{\text{eff}}$ is probed in Section~4.6. A representativity study in the sense of Hill~[4], i.e., a statistical assessment over microstructure realizations, sample sizes and loading conditions, is outside the scope of this work.}{Independence from the initial microscopic crack configuration and adequate representation of fracture mechanisms can be ensured by selecting a sufficiently large representative region that contains an appropriate number of microstructural features.
This presumes a sufficiently homogeneous spatial distribution of microstructural features, analogous to assumptions made for representative volume elements in micromechanics.
In the present work, this requirement is not fully verified: only one prescribed macroscopic crack-growth direction and one initial crack location are considered, and no convergence study with respect to the size of the resolved porous region is performed.}
%@generated-end
\end{relatedchanges}

\begin{reviewercomment}
12. There is still text between section and subsection (3. vs. 3.1). I complained about this issue the last time already.
\end{reviewercomment}
\begin{response}
We apologize. The sentence between the heading of Section~3 and that of Section~3.1 was moved into Section~3.1. Sections~4 and 5 start directly with their first subsection.
\end{response}
\begin{relatedchanges}
\relhead{Sections~3 and 3.1}
%@quote cmd:This section introduces the phase-field models and numerical ingredients
%@quote cmd:This section introduces the phase-field models and the numerical ingredients
%@generated-begin
\item \msloc{Section~3}
\deleted{This section introduces the phase-field models and numerical ingredients used in the subsequent numerical experiments.}
\item \msloc{Section~3.1}
\added{This section introduces the phase-field models and the numerical ingredients used in the subsequent numerical experiments.}
%@generated-end
\end{relatedchanges}

\begin{reviewercomment}
13. comma missing at the end of eq. (13)
\end{reviewercomment}
\begin{response}
Done.
\end{response}
\begin{relatedchanges}
\relhead{Section~3.1, eq.~(13)}
%@quote env:f\,\dot{\gamma} = 0\changed
%@generated-begin
\item \msloc{Section~3.1}
\begin{equation*}
    f(\boldsymbol{\sigma},\alpha) \le 0, 
    \quad \dot{\gamma} \ge 0, 
    \quad f\,\dot{\gamma} = 0\changed{,}
\end{equation*}
%@generated-end
\end{relatedchanges}

\begin{reviewercomment}
14. ``serving as a residual stiffness parameter to avoid ill-conditioned tangents if the problem is solved using a Newton
scheme.'' Well, this has nothing to do with the numerics. Rather, the model itself is severely ill-posed without
positive eta.
\end{reviewercomment}
\begin{response}
We agree. The remark on the numerics was removed, and the text now states that a positive $\eta$ keeps the stiffness positive definite in fully fractured regions, so that the equilibrium problem remains well-posed, and that $\eta$ is a part of the model.
\end{response}
\begin{relatedchanges}
\relhead{Section~3.1, after eq.~(16)}
%@quote cmd:keeps the stiffness positive definite
%@generated-begin
\item \msloc{Section~3.1}
\replaced{a small residual stiffness. A positive value of $\eta$ keeps the stiffness positive definite in fully fractured regions, $s = 0$, so that the equilibrium problem remains well-posed there. It is therefore a part of the model and not a device of the numerical solution procedure.}{serving as a residual stiffness parameter to avoid ill-conditioned tangents if the problem is solved using a Newton scheme.}
%@generated-end
\end{relatedchanges}

\begin{reviewercomment}
15. ``no pointwise convergence statement for s is intended here,'' I understand what is meant, but the sentence itself
does not make sense in its current form.
\end{reviewercomment}
\begin{response}
The sentence was rewritten.
\end{response}
\begin{relatedchanges}
\relhead{Section~3.1, after eq.~(18)}
%@quote cmd:No pointwise convergence of $s$ is implied
%@generated-begin
\item \msloc{Section~3.1}
\replaced{In this sense, the regularized fields $(\boldsymbol{u},s)$ approximate the displacement field and the crack set of the corresponding sharp-crack formulation. The convergence is that of the energy functionals and of their minimizers, and the region of low values of $s$ concentrates on the crack set as $\epsilon \to 0$. No pointwise convergence of $s$ is implied, see Kuhn~[39] and Chambolle~[40].}{In this sense, the regularized fields $(\boldsymbol{u},s)$ approximate the displacement field and crack set of the corresponding sharp-crack formulation, while the diffuse crack region represented by low values of $s$ approaches the crack set; no pointwise convergence statement for $s$ is intended here, see Kuhn~[39] and Chambolle~[40].}
%@generated-end
\end{relatedchanges}

\begin{reviewercomment}
16. ``The mobility parameter M arises from the introduction of a viscous regularization and facilitates the numerical
solution of fracture problems under quasi-static loading conditions.'' This is reverse logic! The parameter M is
introduced to write down the phase-field evolution eq. (19), not the other way around.
\end{reviewercomment}
\begin{response}
We agree. The text now introduces $M$ as the kinetic coefficient of the evolution equation~(19), states that large values of $M$ approximate a rate-independent evolution, and refers to the energy partition of Section~\mycomment{TODO}{4.3}, where the energy dissipated by the viscous evolution of $s$ is quantified for the numerical experiments.
\end{response}
\begin{relatedchanges}
\relhead{Section~3.1, after eq.~(19)}
%@quote cmd:kinetic coefficient of the evolution equation
%@generated-begin
\item \msloc{Section~3.1}
\replaced{The parameter $M$ is the kinetic coefficient of the evolution equation~(19). It sets the time scale on which the phase field relaxes towards a stationary point of the energy, and for large values of $M$ the evolution of $s$ approximates a rate-independent process. The energy dissipated by this viscous evolution, $\int_\Omega \dot{s}^2/M\,\mathrm{d}V$ per unit time, is part of the energy balance of the model and is quantified for the numerical experiments in Section~4.3.}{The mobility parameter $M$ arises from the introduction of a viscous regularization and facilitates the numerical solution of fracture problems under quasi-static loading conditions.}
%@generated-end
\end{relatedchanges}

\begin{reviewercomment}
17. Eq. (21) appears late. The isotropy assumption on the stiffness should be stated earlier. Moreover, the shear and
compression modulus should be introduced.
\end{reviewercomment}
\begin{response}
The isotropy of the bulk material is now stated where the elastic energy is introduced, eq.~(15), together with the bulk modulus $K$, the shear modulus $\mu$ and their relation to the Lam\'e parameters. Eq.~(21) is now introduced as the consequence of eq.~(15) for the isotropic stiffness.
\end{response}
\begin{relatedchanges}
\relhead{Section~3.1, eqs.~(15) and (21)}
%@quote cmd:characterized by the bulk modulus
%@quote cmd:follows from eq.~(\ref{eq:psi_el}) for the isotropic stiffness
%@generated-begin
\item \msloc{Section~3.1}
\added{ of the undamaged material. Throughout this work, the bulk material is assumed to be isotropic, so that $\boldsymbol{\mathbb{C}}$ is characterized by the bulk modulus $K$ and the shear modulus $\mu$, or equivalently by the Lam\'e parameters $\lambda = K - \nicefrac{2}{3}\mu$ and $\mu$. The stiffness}
\item \msloc{Section~3.1}
\replaced{follows from eq.~(15) for the isotropic stiffness as}{is computed under the assumption of isotropic linear elasticity as}
%@generated-end
\end{relatedchanges}

\begin{reviewercomment}
18. ``Gateaux derivative'' $\to$ ``derivative''
\end{reviewercomment}
\begin{response}
Done, see the quote under the next point.
\end{response}

\begin{reviewercomment}
19. section 3.3: This is the proper order: first discretizaion, then assembling, then solver.
\end{reviewercomment}
\begin{response}
Section~3.3 was reordered accordingly: spatial and temporal discretization first, then the assembly of the monolithic system, then the solvers and the time step control.
\end{response}
\begin{relatedchanges}
\relhead{Section~3.3}
%@quote cmd:The spatial discretization uses first-order Lagrange finite elements on triangular meshes in two dimensions and tetrahedral meshes in three dimensions. Along the expected crack path, the mesh is locally refined such that the element edge length $\Delta h$ is at most half
%@quote cmd:The weak forms of the coupled boundary-value problems are implemented in FEniCSx using \texttt{dolfinx} version 0.7.3
%@generated-begin
\item \msloc{Section~3.3}
\added{The spatial discretization uses first-order Lagrange finite elements on triangular meshes in two dimensions and tetrahedral meshes in three dimensions. Along the expected crack path, the mesh is locally refined such that the element edge length $\Delta h$ is at most half the regularization length $\epsilon$, and it is usually chosen much smaller. The specific values of $\Delta h$ and $\epsilon$ used in the individual simulations are stated together with the corresponding numerical setups below. The evolution equation for the phase field, eq.~(19), is discretized in pseudo-time by an implicit backward Euler scheme.}
\item \msloc{Section~3.3}
\added{The weak forms of the coupled boundary-value problems are implemented in FEniCSx using \texttt{dolfinx} version 0.7.3~[42]. The displacement field, the phase field, and, where applicable, the plastic internal variables form one monolithic nonlinear system per time step. The nonlinear residual and its derivative are assembled by FEniCSx, and the resulting system is solved with the PETSc-based Newton solver provided by FEniCSx. The linear subproblems arising within Newton's method are solved with the default PETSc Krylov solver GMRES and the default PETSc ILU preconditioner used by this setup. The time step is adapted automatically during the simulation. The maximum number of Newton iterations per time step is set to eight. If a step does not converge within this limit, the time step is reduced by a factor of two and the step is repeated. If convergence is reached in fewer than four Newton iterations, the next time step is increased by a factor of two, subject to the prescribed upper bound for the respective simulation.}
%@generated-end
\end{relatedchanges}

\begin{reviewercomment}
20. ``satisfies $\Delta h \le \epsilon/2$''; inequalities are no words (see previous review)
\end{reviewercomment}
\begin{response}
Reworded in words (``at most half the regularization length''), see the quote under point~19.
\end{response}

\begin{reviewercomment}
21. ``In the following chapters,''; there are no chapters, only sections or paragraphs
\end{reviewercomment}
\begin{response}
Done.
\end{response}
\begin{relatedchanges}
\relhead{Section~3.4}
%@quote cmd:sections}{chapters}
%@generated-begin
\item \msloc{Section~3.4}
\replaced{sections}{chapters}
%@generated-end
\end{relatedchanges}

\begin{reviewercomment}
22. ``crack nucleation in previously undamaged material can also be predicted using phase-field models.'' This is a
strong statement. It is true that one can apply phase-field fracture to pristine setups. Whether the results are
\emph{predictive}, is another matter entirely.
\end{reviewercomment}
\begin{response}
We agree. ``Predicted'' was replaced by ``modelled'', and the sentence was qualified: whether the nucleation loads are quantitatively predictive is a separate question, and nucleation is used in this work as a model mechanism with analytically known characteristic stress and energy.
\end{response}
\begin{relatedchanges}
\relhead{Section~3.6}
%@quote cmd:modelled with}{predicted using}
%@quote cmd:Whether the nucleation loads obtained in this way
%@generated-begin
\item \msloc{Section~3.6}
\replaced{modelled with}{predicted using}
\item \msloc{Section~3.6}
\added{ Whether the nucleation loads obtained in this way are quantitatively predictive for a given material is a separate question, which depends on the regularization length and on the degradation function, see Tann\'e et al.~[43]. In the present work, nucleation is used as a model mechanism whose characteristic stress and energy are known analytically.}
%@generated-end
\end{relatedchanges}

\begin{reviewercomment}
23. ``from (21), cf. [28].'' (see previous review)
\end{reviewercomment}
\begin{response}
Corrected.
\end{response}
\begin{relatedchanges}
\relhead{Section~3.6}
%@quote cmd:, see Noll et al.}{, cf.}
%@generated-begin
\item \msloc{Section~3.6}
\replaced{, see Noll et al.}{, cf.}
%@generated-end
\end{relatedchanges}

\begin{reviewercomment}
24. ``Strictly speaking, these relations hold only in one dimension, but they provide good estimates for two- and
three-dimensional scenarios,'' Evidence, references, or quantitative error estimates should be provided to support this
claim.
\end{reviewercomment}
\begin{response}
The statement was too broad. It now refers to the evidence that exists for two and three dimensions (Kuhn 2013 for the quadratic and Schl\"uter 2018 for the cubic degradation function: nucleation once a stress close to $\sigma_c$ is exceeded in a region of the size of several $\epsilon$), and the new simulations record the stress state in the ligament ahead of the trapped crack tip immediately before re-nucleation, which is compared with the one-dimensional values in Section~\mycomment{TODO}{4.4 (study S1 outputs; numbers to be inserted)}.
\end{response}
\begin{relatedchanges}
\relhead{Section~3.6}
%@quote cmd:Strictly speaking, these relations hold only in one dimension. In two and three dimensions
%@generated-begin
\item \msloc{Section~3.6}
\replaced{Strictly speaking, these relations hold only in one dimension. In two and three dimensions, cracks nucleate once a stress level close to $\sigma_c$ is exceeded in a region whose size is of the order of several multiples of $\epsilon$, as shown for the quadratic degradation function by Kuhn~[39] and for the cubic degradation function by Schl\"uter~[44]. For the porous specimens studied in this work, the stress state in the ligament ahead of a trapped crack tip immediately before re-nucleation is recorded and compared with the one-dimensional values in Section~4.4.}{Strictly speaking, these relations hold only in one dimension, but they provide good estimates for two- and three-dimensional scenarios, where cracks tend to form spontaneously once a stress level close to $\sigma_c$ is exceeded in a sufficiently large region—typically on the order of several multiples of $\epsilon$, cf.~[39].}
%@generated-end
\end{relatedchanges}

\begin{reviewercomment}
25. The plotted figures have low quality, e.g., Fig. 9. Also, the axis labels should start with an upper case.
\end{reviewercomment}
\begin{response}\mycomment{TODO}{all result figures regenerated from the new data as vector graphics (pgf), axis labels capitalised.}\end{response}

\begin{reviewercomment}
26. ``compared the results against'' $\to$ ``compared the results to/with''
\end{reviewercomment}
\begin{response}
Done.
\end{response}
\begin{relatedchanges}
\relhead{Conclusion}
%@quote cmd:with}{against}
%@generated-begin
\item \msloc{Conclusion}
\replaced{with}{against}
%@generated-end
\end{relatedchanges}

%==================================================================================================
\section*{Reviewer 2}

\begin{reviewercomment}
Revisions improved the manuscript.
\end{reviewercomment}

\begin{response}
We thank the reviewer for the positive assessment.
\end{response}

%==================================================================================================
\section*{Reviewer 3}

\begin{reviewercomment}
I thank the authors for their careful revisions. The scope and positioning of the manuscript are now clearer, and I
consider that it can be accepted after the following minor points are addressed.
\end{reviewercomment}

\begin{response}
We thank the reviewer for the positive assessment and for the three comments, which we address below.
\end{response}

\begin{reviewercomment}
Comment 1: The main remaining weakness is the limited discussion of the mechanisms behind the numerical observations.
The authors describe their simulations as numerical experiments, but experimental papers usually go beyond
observations. The authors have the chance to have at hand much more observables (including energetic ones) than
experimentally available.

In particular, I think that the impact of plastic dissipation on the nucleation (small $w_s/L$ regime) and propagation
(large $w_s/L$ regime) in porous materials need to be discussed more thoroughly. I do not request additional
simulations. While it is quite intuitive to understand why bulk dissipation inhibits crack growth (large $w_s/L$
regime), it is not clear why ``the ductile material law leads to a reduction in effective crack resistance below''. A
focused discussion based on the stress distribution in the immediate vicinity of the crack tip/void wall where crack
renucleate, and an energy partitioning, based on Eq. (14) would be sufficient.
\end{reviewercomment}

\begin{response}
\mycomment{TODO}{same material as Reviewer 1, major point 3: energy partition based on eq.~(14) (now with the exact balance incl.\ the viscous term), stress/plastic-zone fields in the ligament immediately before re-nucleation for $w_s = 0.375L$ and $1.0L$ (vM vs.\ Eq$\sigma^*$), explanation of the small-$w_s$ regime (plastic zone spans the ligament, strain-driven nucleation at lower stress, the plastic work spent before nucleation is not available for the jump) vs.\ the large-$w_s$ regime (plastic dissipation ahead of the trapped tip delays re-nucleation).}
\end{response}

\begin{relatedchanges}
\relhead{\mycomment{TODO}{quotes as for R1 major point 3}}
\end{relatedchanges}

\begin{reviewercomment}
Comment 2: When the authors state ``Whether such a fracture toughness exists for heterogeneous materials and under which
circumstances both concepts might coincide remains an open question of research'', they are not entirely correct.
Braides et al. (1996, doi:10.1007/BF02198476) and more recently Cagnetti et al. (2019, doi:10.1007/s00205-019-01372-x)
showed that, under some assumptions, the fracture energy, interpreted as a surface energy, could be homogenized. This
is discussed in details by Michel and Suquet (2022, cited by the authors). As pointed by the authors, this is ``not
established in general'' but results exist and should not be attributed to Schneider only.
\end{reviewercomment}

\begin{response}
We thank the reviewer for the references. The statement is corrected: the homogenization of a surface energy of
Griffith type is established for periodic and, under assumptions, for stochastic microstructures (Braides, Defranceschi
and Vitali 1996; Cagnetti, Dal Maso, Scardia and Zeppieri 2019), and the variational definition used by Schneider builds
on these results, as discussed by Michel and Suquet. The open question addressed in our work is a different one, namely
whether the maximum driving force recorded in a crack-growth experiment with an elasto-plastic bulk material and
unstable re-nucleation can be identified with such a homogenized surface energy. The references of Braides et al.\ (1996)
and Cagnetti et al.\ (2019) were added to the bibliography, together with Francfort and Marigo (1998) for the variational
formulation whose surface energy is homogenized.
\end{response}

\begin{relatedchanges}
\relhead{Introduction}
%@quote cmd:For elastic and plastic properties, the step from the microstructure
%@quote cmd:Rigorous results exist for the homogenization of surface energies
%@generated-begin
\item \msloc{Introduction}
\added{For elastic and plastic properties, the step from the microstructure to a macroscopic material model is provided by homogenization. Its central concept is the representative volume element, defined by Hill~[4] as the smallest sample of a heterogeneous material that can be treated as a statistically uniform continuum. The effective properties of such a sample are independent of the boundary conditions applied to it, see also Nemat-Nasser and Hori~[5]. Fracture does not fit into this framework in a direct way. A crack is a localized and evolving discontinuity, its growth is a non-stationary process that may be unstable, and the separation of length scales that underlies the notion of a representative volume element is lost once the crack interacts with individual microstructural features. For this reason, the term \textsl{fracture toughness} is used in this work only for a macroscopic material property in the sense of eq.~(1), whereas the value obtained from a numerical experiment under specific conditions, i.e., the maximum energetic driving force required to force crack advance through the considered microstructure, is called the \textsl{effective crack resistance}. Three classes of approaches to the latter quantity can be distinguished in the literature.}
\item \msloc{Introduction}
\added{Rigorous results exist for the homogenization of surface energies of Griffith type. Braides, Defranceschi and Vitali~[6] proved that free-discontinuity functionals with a periodically oscillating surface energy, as they appear in the variational formulation of brittle fracture by Francfort and Marigo~[7] and Bourdin et al.~[8], $\Gamma$-converge to a functional with a homogenized and, in general, anisotropic surface energy. Cagnetti, Dal Maso, Scardia and Zeppieri~[9] extended this result to stationary random microstructures. Schneider~[10] and Ernesti and Schneider~[11, 12] use the resulting cell formula as a variational definition of an effective crack energy and compute it with FFT-based solvers, also for anisotropic fracture resistance. Michel and Suquet~[13] assess the merits and limitations of this definition. They relate it to limit-load problems and show that it is a lower bound of a more physically motivated effective fracture energy, namely the energy dissipated along the actual crack path through a unit cell. The homogenized surface energy is the energy of the minimal crack surface weighted by the local toughness. It does not account for the elastic interaction of the crack with the microstructure, for crack trapping and the unstable crack jumps that follow re-nucleation, or for inelastic bulk dissipation, which are the mechanisms that govern toughening in many materials, see the discussion by Lebihain~[14].}
%@generated-end
\end{relatedchanges}

\begin{reviewercomment}
Comment 3: The definition of the different models vM, Eq$\sigma$ and Eq$J_c$ are now clear, and we understand their
purpose. However (and in my opinion), to mirror the link between vM and Eq$\sigma$ and minimize finite-size effects,
Eq$J_c$ could have an evaluated $J_c^{\mathrm{eff}}$ equal to that of vM, not merely setting $J_c^0$. Repeating the
porous simulations is not necessary and a brief discussion of the expected effect would be enough.
\end{reviewercomment}

\begin{response}
We agree that a linear elastic reference set whose numerical crack resistance $J_c^{\mathrm{num}}$ in the homogeneous
strip equals that of the von~Mises set would be a third natural choice. Its effect can be predicted without
simulations: for the rate-independent linear elastic phase-field model the only energy scale is $J_c$, so that all
energies, and in particular $J_{x,c}^{\text{eff}}$, scale linearly with $J_c$ while the crack pattern is unchanged (the
displacement field scales with $\sqrt{J_c}$). The two linear elastic sets of the manuscript confirm this: normalised by
their own $J_c^{\mathrm{num}}$, \EqJcLabel and \EqSigmaStarLabel give the same curve $J_{x,c}^{\text{eff}}(w_s)$
\mycomment{TODO}{insert the numbers from the new S1 runs; v2.0 data: both $\approx 1.9$--$2.0\,J_c^{\mathrm{num}}$ in
the saturation regime}. A linear elastic set with $J_c^{\mathrm{num}} = J_c^{\mathrm{num,pl}}$ would therefore reach
$J_{x,c}^{\text{eff}} \approx 1.9\,J_c^{\mathrm{num,pl}}$, i.e.\ it would be \emph{tougher} than the von~Mises set
($\approx 1.25\,J_c^{\mathrm{num,pl}}$ \mycomment{TODO}{new value}) at large wall widths. This strengthens the
conclusion that the von~Mises material does not simply behave like a linear elastic material with an increased crack
resistance; the additional dissipation during re-nucleation is spent differently. A paragraph with this argument was
added. \mycomment{TODO}{quote}
\end{response}

\begin{relatedchanges}
\relhead{\mycomment{TODO}{quote: new paragraph at the end of Section~4.2}}
\end{relatedchanges}

%==================================================================================================
\section*{Additional corrections by the authors}

\begin{response}
While preparing the new simulations we audited the simulation code against the manuscript and corrected the following
inconsistencies. None of them changes the conclusions of the paper, but the quantitative results were recomputed with
consistent parameters.
\begin{enumerate}
\item The fracture energy density, eq.~(18), was stated with the regularisation $(1-s)^2/(2\epsilon) + 2\epsilon|\nabla s|^2$,
whereas the implemented crack surface density is $(1-s)^2/(4\epsilon) + \epsilon|\nabla s|^2$, which integrates to the
crack length for the one-dimensional profile and is consistent with the evolution equation~(19) and the nucleation
stress~(26) given in the manuscript. Eq.~(18) was corrected.
\item The Kolosov constant in the surfing boundary condition, eq.~(2), is $\kappa = 3 - 4\nu_{\mathrm{bc}}$ for the plane
strain setting used throughout; the plane stress expression was printed by mistake.
\item The auxiliary parameters of the K-field in the porous experiments are the effective elastic constants of the
surrounding material, $E_{\mathrm{bc}} = E_0^{\mathrm{eff}}$ and $\nu_{\mathrm{bc}} = \nu_0^{\mathrm{eff}}$ (Sections~4.2
and 5.1); the values $2.5\mu_0$ and $0.25$ hold only for the homogeneous strip of Section~3.5.
\item The slope parameter of the cubic degradation function was $\beta = 0.1$ in the porous and the homogeneous-strip
simulations of the previous version but $0.2$ in the one-dimensional bar; the stated value $\beta = 0.2$ is now used in
all simulations. Likewise, the residual stiffness of the three-dimensional simulations was $\eta = 10^{-3}$ and is now
$10^{-5}$ as stated, and the yield stress is assigned to the porous region only, so that the surrounding material is
linear elastic as described. All results in Sections~3.5, 4 and 5 were recomputed accordingly \mycomment{TODO}{(list the
numbers that changed: $J_c^{\mathrm{num}}$, saturation values, \dots)}.
\item The FEniCSx version is 0.7.3 (not 0.7.4).
\item The one-dimensional bar of Section~3.6 was recomputed with the parameters now used throughout ($\beta = 0.2$, $\eta =
10^{-5}$, $\epsilon = 0.1L$, $M = 1000$, $\Delta h = 0.02L$). The fracture stress of the von~Mises bar is $\sigma^* = 1.28\sigma_y$
(previously $1.31\sigma_y$, a value belonging to $\beta = 0.1$); the linear elastic bar with $J_c = J_c^0$ is unchanged at
$1.72\sigma_y$. Since the fracture stress of the linear elastic model is proportional to $\sqrt{J_c}$, the parameter set
\EqSigmaStarLabel{} now uses $J_c = 0.5523J_c^0$ (previously $0.5679J_c^0$, which gave $1.30\sigma_y$) and has exactly the
fracture stress of the von~Mises bar. The fracture strains and dissipated energies were updated accordingly
($\varepsilon^* = 1.17,\ 1.07,\ 0.80$ and $\Pi^* = 2.12,\ 2.19,\ 1.23\,J_c^0L$ for \vMLabel, \EqJcLabel, \EqSigmaStarLabel),
the fracture strain is now defined explicitly as the strain at which the stress drops to zero, and the width of the
weakened zone that triggers nucleation is stated as implemented ($|x| \le 0.15L$). The finite-element peaks agree with a
rate-independent homogeneous solution of the one-dimensional problem to within $0.1\,\%$, and halving the loading rate or
increasing the mobility tenfold changes them by less than $0.5\,\%$. Figures~5 and~6 were regenerated.
\item The effective elastic constants of the two-dimensional unit cells reported in Section~4.1 of the previous versions
were not homogenized values but the Voigt bound $(1-f)\,\mathbb{C}_0$ of the porous cell (porosity $f$): the linear
displacement boundary condition of the homogenization step had inadvertently also been applied to the pore surface,
which makes the displacement field affine. The homogenization was corrected (linear displacements on the outer edges of
the unit cell only, traction-free pore surface); the corrected effective stiffness is lower and of cubic symmetry, see
the response to Reviewer~1, point~4. The three-dimensional unit cells and the homogeneous strip were not affected.
For $w_s = 0.375L$, $1.0L$ and $4.0L$ the isotropic approximation of the effective Young's modulus changes from
$E_0^{\mathrm{eff}} = 1.46,\ 2.01,\ 2.42\,\mu_0$ (previous versions, $\nu_0^{\mathrm{eff}} = 0.25$ in all cases) to
$0.90,\ 1.52,\ 2.29\,\mu_0$ with $\nu_0^{\mathrm{eff}} = 0.195,\ 0.228,\ 0.249$; the Zener ratios of the corrected
stiffness are $A = 0.89,\ 0.96,\ 1.00$ (all values for $\lambda_0 = \mu_0$, complete table in Section~4.1).
\end{enumerate}
\end{response}

\begin{relatedchanges}
\relhead{Corrections of equations and parameters}
%@quote env:eq:psi_fr
%@quote cmd:The surfing displacements with $u_z = 0$ are prescribed on the entire outer boundary
%@quote env:eq:surfing-bcs
%@quote cmd:$J_c=0.5523J_c^0$
%@quote cmd:equals}{is close to}
%@quote cmd:$x=-7.5\Delta h$
%@quote cmd:1.31\sigma_y
%@quote cmd:defined as the strain at which the stress drops to zero
%@quote cmd:1.07\sqrt
%@quote cmd:1.17\sqrt
%@quote cmd:$J_c = 0.5523J_c^0$
%@quote cmd:proportional to $\sqrt{J_c}$
%@quote cmd:the same fracture stress of
%@quote cmd:0.80\sqrt
%@quote cmd:2.19J_c^0 L
%@quote cmd:2.12J_c^0 L
%@quote cmd:1.23J_c^0 L
%@quote cmd:Kolosov constant
%@quote cmd:The weak forms of the coupled boundary-value problems are implemented in FEniCSx using \texttt{dolfinx} version 0.7.3
%@quote cmd:effective elastic constants of the surrounding material, $E_{\text{bc}} = E_0^{\text{eff}}$
%@quote cmd:as in the two-dimensional case, the effective elastic constants
%@generated-begin
\item \msloc{Section~2, Eq.~(2)}
\begin{align*}
    \begin{split}
        u_x &= \dfrac{K_I^{\text{bc}}}{2\mu_{\text{bc}}}  \sqrt{\dfrac{r}{2\pi}}\left(\kappa - \cos{\varphi}\right)\cos{\dfrac{\varphi}{2}},\\
        u_y &= \dfrac{K_I^{\text{bc}}}{2\mu_{\text{bc}}}  \sqrt{\dfrac{r}{2\pi}}\left(\kappa - \cos{\varphi}\right)\sin{\dfrac{\varphi}{2}},\\
        \text{where} \quad \kappa &= \changed{3 - 4\nu_{\text{bc}}} \quad \text{and} \quad \mu_{\text{bc}} = \dfrac{E_{\text{bc}}}{2\left(1 + \nu_{\text{bc}}\right)}.
    \end{split}
\end{align*}
\item \msloc{Section~2}
\added{, and $\kappa$ is the Kolosov constant for plane strain}
\item \msloc{Section~3.1, Eq.~(18)}
\begin{equation*}
    \psi_{\rm{fr}} = J_c
    \underbrace{\left(
    \changed{\dfrac{(1-s)^2}{4\epsilon} + \epsilon\|\nabla{s}\|^2}
    \right)}_{\bar{A}},
\end{equation*}
\item \msloc{Section~3.3}
\added{The weak forms of the coupled boundary-value problems are implemented in FEniCSx using \texttt{dolfinx} version 0.7.3~[42]. The displacement field, the phase field, and, where applicable, the plastic internal variables form one monolithic nonlinear system per time step. The nonlinear residual and its derivative are assembled by FEniCSx, and the resulting system is solved with the PETSc-based Newton solver provided by FEniCSx. The linear subproblems arising within Newton's method are solved with the default PETSc Krylov solver GMRES and the default PETSc ILU preconditioner used by this setup. The time step is adapted automatically during the simulation. The maximum number of Newton iterations per time step is set to eight. If a step does not converge within this limit, the time step is reduced by a factor of two and the step is repeated. If convergence is reached in fewer than four Newton iterations, the next time step is increased by a factor of two, subject to the prescribed upper bound for the respective simulation.}
\item \msloc{Section~3.4}
\replaced{$J_c=0.5523J_c^0$}{$J_c=0.5679J_c^0$}
\item \msloc{Section~3.4}
\replaced{equals}{is close to}
\item \msloc{Section~3.6}
\replaced{$x=-7.5\Delta h$ to $x=7.5\Delta h$, i.e.\ $|x| \le 0.15L$}{$x=-3\Delta h$ to $x=3\Delta h$}
\item \msloc{Section~3.6}
\replaced{$\sigma^* = 1.28\sigma_y$}{$\sigma^* = 1.31\sigma_y$}
\item \msloc{Section~3.6}
\added{ $\varepsilon^*$, defined as the strain at which the stress drops to zero,}
\item \msloc{Section~3.6}
\replaced{$\varepsilon^* \approx 1.07\sqrt{\nicefrac{J_c^0}{\mu_0 L}}$}{$\varepsilon^* \approx 1.08\sqrt{\nicefrac{J_c^0}{\mu_0 L}}$}
\item \msloc{Section~3.6}
\replaced{$\varepsilon^* \approx 1.17\sqrt{\nicefrac{J_c^0}{\mu_0 L}}$}{$\varepsilon^* \approx 1.18\sqrt{\nicefrac{J_c^0}{\mu_0 L}}$}
\item \msloc{Section~3.6}
\replaced{$J_c = 0.5523J_c^0$}{$J_c = 0.5679J_c^0$}
\item \msloc{Section~3.6}
\added{, since the fracture stress of the linear elastic model is proportional to $\sqrt{J_c}$}
\item \msloc{Section~3.6}
\replaced{the same fracture stress of $\sigma^* = 1.28\sigma_y$ as for the von~Mises model}{an almost equivalent fracture stress of $\sigma^* = 1.30\sigma_y$ compared to the von~Mises model}
\item \msloc{Section~3.6}
\replaced{$\varepsilon^* \approx 0.80\sqrt{\nicefrac{J_c^0}{\mu_0 L}}$}{$\varepsilon^* \approx 0.82\sqrt{\nicefrac{J_c^0}{\mu_0 L}}$}
\item \msloc{Section~3.6}
\replaced{$\Pi^* \approx 2.19J_c^0 L$}{$\Pi^* \approx 2.22J_c^0 L$}
\item \msloc{Section~3.6}
\replaced{$\Pi^* \approx 2.12J_c^0 L$}{$\Pi^* \approx 2.16J_c^0 L$}
\item \msloc{Section~3.6}
\replaced{$\Pi^* \approx 1.23J_c^0 L$}{$\Pi^* \approx 1.29J_c^0 L$}
\item \msloc{Section~4.2}
\replaced{ and the effective elastic constants of the surrounding material, $E_{\text{bc}} = E_0^{\text{eff}}$ and $\nu_{\text{bc}} = \nu_0^{\text{eff}}$, as the auxiliary parameters of the prescribed K-field~(2)}{, $E_{\text{bc}} = 2.5 \mu_0$, and $\nu_{\text{bc}} = 0.25$}
\item \msloc{Section~5.1}
\replaced{ and, as in the two-dimensional case, the effective elastic constants of the surrounding material as the auxiliary parameters $E_{\text{bc}}$ and $\nu_{\text{bc}}$}{, $E_{\text{bc}}=2.5\mu_0$, and $\nu_{\text{bc}}=0.25$}
\item \msloc{Section~5.1}
\replaced{The surfing displacements with $u_z = 0$ are prescribed on the entire outer boundary, including the planes $z=z_{\rm min}$ and $z=z_{\rm max}$, which mimics}{Additionally, the displacement $u_z$ is constrained to zero on the planes $z=z_{\rm min}$ and
$z=z_{\rm max}$ to mimic}
%@generated-end
\end{relatedchanges}

\end{document}
